\def\level{BTS}  % Classe à droite
\def\course{CIEL 2}
\def\chaptername{Lois continues} % Titre au centre
\def\docpart{Loi exponentielle~--~Exercices} % Partie du cours
\def\docname{C2253}        % Identifiant à gauche

\pagestyle{cours_FP}
\makeheaderBTS{} % Affiche le bandeau

\begin{multicols}{2}
\begin{exercice_cours}[\quad]{}
La durée de vie d'un composant électronique suit une loi exponentielle de paramètre $\lambda=0,225$.

\begin{enumerate}
	\item Quelle est la probabilité qu'un composant dure moins de 8 ans ? Plus de 8 ans ?
	\item Quelle est la probabilité qu'un composant dure plus de 8 ans sachant qu'il a déjà duré plus de 3 ans ?
\end{enumerate}


\end{exercice_cours}
\begin{exercice_cours}[\quad]{}
La durée d'attente à une caisse de supermarché, exprimée en minutes, est une variable aléatoire qui suit une loi exponentielle de paramètre $\lambda=0,04$.

\begin{enumerate}
	\item Quelle est la probabilité qu'un client attende moins de 5 minutes ?
	\item Quelle est la probabilité qu'un client attende plus de 10 minutes ?
\end{enumerate}

\end{exercice_cours}
\begin{exercice_cours}[\quad]{}
La durée de vie, exprimée en heures, d'un téléphone portable est une variable aléatoire X qui suit une loi exponentielle de paramètre $\lambda=0,00026$.

\begin{enumerate}
	\item Calculer P(X$\leq 1000$) et P(X$\geq 1000$).
	\item Sachant que l'événement (X$\geq 1000$) est réalisé, calculer la probabilité de l'événement (X$\geq 2000$).
\end{enumerate}

\end{exercice_cours}
\begin{exercice_cours}[\quad]{}
On considère une production de lampes pour vidéoprojecteurs. On admet que la variable aléatoire T qui, à toute lampe choisie au hasard dans la production, associe sa durée de vie, suit une loi exponentielle de paramètre $\lambda$.

\begin{enumerate}
	\item Le constructeur annonce que la durée de vie moyenne d'une lampe est de $2000$ heures. En déduire la valeur de $\lambda$.
	\item Quelle est la probabilité que la durée de vie d'une lampe soit supérieure à $4000$ heures ?
	\item Déterminer le réel $t$ tel que P(T$\leq t$) = 0,7.
	\item Sachant qu'une lampe a fonctionné plus de $2000$ heures, quelle est la probabilité pour qu'elle tombe en panne avant $4000$ heures ?
\end{enumerate}


\end{exercice_cours}
\begin{exercice_cours}[\quad]{}
La durée de vie d'un robot, exprimée en années, est une vaiable aléatoire qui suit une loi exponentielle de paramètre $\lambda=0,2$.

\begin{enumerate}
	\item À quel instant $t$, à un mois près, la probabilité qu'un robot tombe en panne pour la première fois est-elle de $0,5$ ?
	\item Montrer que la probabilité qu'un robot n'ait pas eu de panne au cours des deux premiètres années est $e^{-0,4}$.
\end{enumerate}

\end{exercice_cours}
\begin{exercice_cours}[\quad]{}
La durée de vie d'un composant électronique jusqu'à ce que survienne la première panne, exprimée en années, est une variable aléatoire X qui suit une loi exponentielle de paramètre $\ell$ ($\ell > 0$).

\begin{enumerate}
	\item Déterminer $\ell$, arrondi à $10^{-4}$ près, pour que la probabilité P(X $> 8$) soit égale à $0,2$.
	
	Dans la suite de l'exercice, on prendra $\ell = 0,2$.
	%\item A quel instant $t$, à un mois près, la probabilité qu'un composant électronique tombe en panne pour la première fois est-elle de $0,5$ ?
	%\item Montrer que la probabilité qu'un composant électronique n'ait pas eu de panne au cours des deux premières années est de $e^{-0,4}$
	\item Sachant qu'un composant électronique n'a pas eu de panne au cours des deux premières années, quelle est à $10^{-2}$ près, la probabilité qu'il soit encore en état de marche au bout de six ans ?
\end{enumerate}

\end{exercice_cours}
\end{multicols}	